\bookchapter{第九章：移動辦公與學習}

\chapterintro{我以前也覺得手機上辦公就是“湊合一下”，直到有一次出差忘帶電腦，硬是用手機完成了標書修改和合同簽署。說實話，手機辦公的瓶頸往往不在手機本身，而在你不熟悉它的能力邊界。}

\booksection{一、文件處理：在掌中編輯與創作}

在手機螢幕上編輯文件，跟電腦上感覺確實不一樣。最大的問題是螢幕小、手指粗、排版容易跑偏。但有幾個技巧能讓體驗好不少：

\begin{itemize}
  \item \textbf{WPS Office}：本土化做得最好，完全相容微軟格式（docx/xlsx/pptx），內建的PDF工具集很豐富。我特別喜歡它的"語音備註"功能——看合同的時候不方便打字，直接對著手機說幾句意見，它就自動轉成文字批註貼在文件裡。另外WPS內建了海量免費模板，做簡歷、做彙報PPT的時候能省很多時間。
  \item \textbf{Microsoft 365}：如果你公司用正版Office，這個就是最佳選擇。跟桌面版的格式相容性幾乎完美，不會出現WPS開啟某些複雜排版會輕微移位的情況。深度整合OneDrive，你在電腦上改到一半的文件，手機上開啟就是最新版本。
  \item \textbf{石墨文件/騰訊文件}：這兩個是"線上協同"的代表。適合多人同時編輯一份文件，每個人改的地方實時顯示。我用石墨文件跟朋友一起做旅行攻略，兩個人同時往表格裡填資訊，效率比一個人弄高多了。
\end{itemize}

\textbf{實用工作流}：我現在的習慣是——重要檔案存OneDrive，日常編輯用WPS，團隊協作開石墨文件。檔案自從上了雲之後，再也沒出現過“忘帶隨身碟”的尷尬。

\booksection{二、筆記應用：從隨手記到知識體系}

記筆記這件事，很多人一上來就想搞一套很複雜的系統，結果堅持不了三天。我的建議是先輕後重——先養成“記”的習慣，再考慮“管理”的問題。

\begin{itemize}
  \item \textbf{隨手記——系統自帶就夠}：iOS備忘錄啟動極快，而且支援掃描文件、手寫、錄音。安卓各品牌的便籤也不差，鎖屏狀態下就能快速啟動。這些工具適合記臨時想法、購物清單、會議要點。關鍵是啟動快，想到了就能記，不用等App載入。
  \item \textbf{稍後讀——收集資訊}：微信的"浮窗"功能其實是個隱藏神器。看到好文章來不及細看，點右上角"浮窗"，它就會縮成一個小圓點貼在螢幕邊緣，隨時可以點開繼續讀。比轉發到"檔案傳輸助手"方便多了。另外，Cubox或Pocket這類稍後讀工具，可以跨平臺收集你在任何地方看到的好內容。
  \item \textbf{搭建體系——Notion和飛書}：當你覺得零散的筆記開始不好找了，就可以上Notion了。這玩意真的很靈活——你可以建一個"讀書筆記資料庫"，自動按書名、作者、打分來組織；也可以搭一個"專案看板"，把工作任務的進度視覺化。飛書的"多維表格"也類似，而且國內訪問速度更快。我用Notion管理我所有的寫作素材、學習筆記和專案進度，已經用了三年多了。
\end{itemize}

\textbf{一個實用的筆記策略}：我自己的習慣是“三件套”——臨時想法用系統便籤（秒開），待讀文章用微信浮窗（順手），需要長期積累的知識放Notion（體系化）。別一上來就想把所有東西塞進一個工具裡，這樣反而容易放棄。

\booksection{三、掃描與OCR：把紙質世界搬進手機}

手機掃描這件事，用好了真的能省掉很多跑列印店的麻煩。我上次辦簽證需要一堆材料的電子版，全都是用手機搞定的。

\begin{itemize}
  \item \textbf{iPhone使用者}：iOS備忘錄自帶的掃描功能已經很好了。開啟備忘錄 點相機圖示 選"掃描文稿"，對準檔案它會自動識別邊框並矯正透視。拍完之後可以調顏色——選"黑白"模式出來的效果最接近掃描器，"灰度"適合有圖片的資料。
  \item \textbf{安卓使用者}：很多品牌的相機自帶"文件掃描"模式。如果找不到，去應用商店下載Microsoft Lens（免費，而且跟OneDrive打通）。
  \item \textbf{OCR文字識別}：拍完的掃描件，如果需要把上面的文字提取出來編輯，就需要OCR了。微信自帶的"提取文字"功能最簡單——長按圖片就能用，識別準確率還不錯。如果大量處理，可以用白描App或者全能掃描王。對於中文識別，百度的OCR引擎效果比較好；WPS的圖片轉文字功能也很準。
  \item \textbf{掃描小技巧}：放在光線充足的地方拍，最好有檯燈從側面打光，避免在文件上留下你的手影。如果檔案折角了，先用手壓平再拍。多頁檔案可以連續掃描後合併成一個PDF，不用一頁頁單獨存。
\end{itemize}

\booksection{四、網盤怎麼選：沒有最好的，只有最適合你的}

選網盤這件事我琢磨了很久，最後發現其實是看你的“主力生態”是什麼。

\begin{itemize}
  \item \textbf{蘋果全家桶 iCloud}：無縫同步到讓你感覺不到它的存在。你在手機上拍的照片，幾分鐘後電腦上就自動有了。而且蘋果的隱私保護做得最嚴格。但缺點也很明顯：跨平臺分享給安卓或Windows使用者比較麻煩，免費只有5GB，擴容要花錢。
  \item \textbf{Windows主力 OneDrive}：如果你用的是Windows電腦+Office 365，OneDrive就是天生一對。Word裡儲存檔案，自動就存到雲上了。1TB空間配合Office訂閱很劃算。手機端App的自動備份照片功能也很好用。
  \item \textbf{國內分享 百度網盤}：說實話，百度網盤的下載限速一直被批評，但它的分享生態確實無人能敵。中文網際網路上大量的學習資料、影視資源都是用百度網盤分享的。如果只是偶爾下載別人分享的檔案，不充會員也能湊合用。如果你是重度的資源收集者，會員該開就開吧，省下的時間值錢多了。
  \item \textbf{阿里雲盤/夸克網盤}：這兩年是後起之秀。阿里雲盤上傳下載不限速（目前是），夸克網盤跟夸克瀏覽器深度繫結，搜尋資料的時候可以直接儲存到網盤，很方便。
\end{itemize}

\textbf{我的組合方案}：iCloud管個人照片和日常檔案（隱私優先），百度網盤存資源和接收分享（生態優先），OneDrive管工作文件（Office整合優先）。三個網盤各司其職，互不干擾。

\booksection{五、一個真實的移動辦公場景}

說個我自己的真實經歷吧。有一次出差到外地，客戶臨時要改合同條款。我的工作流是這樣的：

\begin{enumerate}
  \item 在酒店收到客戶微信發來的修改意見（文字訊息+一張手寫批註的照片）。
  \item 開啟OneDrive找到合同原件，用WPS開啟編輯，參照客戶的微信訊息逐條修改條款。
  \item 改完後用WPS的"輸出為PDF"功能匯出，發到客戶微信群裡確認。
  \item 客戶提了第二輪意見，我用WPS的畫筆工具直接在PDF上圈出修改位置，配合語音批註說清楚改動思路。
  \item 確認無誤後，用WPS的"電子簽名"功能線上簽署，發給對方同樣簽署。
  \item 最後把簽好的合同PDF存檔到百度網盤和iCloud各一份，同時在Notion的工作日誌裡記了一筆。
\end{enumerate}

整個過程從開始到結束大概四十分鐘，全程沒開電腦。說實話，三年前我根本不信手機能幹這些事，但現在確實可以。關鍵是提前熟悉這些工具，別等到需要的時候才手忙腳亂地下載註冊。

\begin{tcbpractice}{來試試：搭建你的移動辦公工作流}
\begin{enumerate}
  \item 選定你的主力文件App（WPS或Microsoft 365），登入並關聯雲盤。
  \item 開啟備忘錄或便籤，用掃描功能拍一份文件，匯出為PDF發給自己。
  \item 選一個網盤App，設定為手機照片自動備份（選"僅WiFi下備份"省流量）。
  \item 在Notion或飛書裡新建一頁，試著搭建一個簡單的"工作日誌"模板。
\end{enumerate}

\end{tcbpractice}
